\section{狭义相对论的两条基本原理}
那么，我们将用什么取代绝对时空观呢？这就是狭义相对论假设的两条基本原理。
\begin{BoxPrinciple}[相对性原理]
    物理定律在所有惯性系中具有相同数学形式，不存在特殊的绝对惯性系。
\end{BoxPrinciple}
\begin{BoxPrinciple}[光速不变原理]
    真空光速$c$是无关惯性系选取的常数，即任何惯性系中测得的真空光速均相同。
\end{BoxPrinciple}
\fancyref{pri:光速不变原理}在数学上应该如何去表达？试想$S(x,y,z,t)$和$S'(x',y',z',t')$是两个不同的惯性系，其中$S'$系相对$S$系以速率$u$沿某个方向运动，那么假设当$S'$和$S$的原点$O'$和$O$重合时，记此时$t'=t=0$时由两者的重合原点发出一个闪光，由于光速与参照系的选定无关，无论在$S'$还是在$S$系中观察，闪光波前均为球面，球心分别为$O'$和$O$，球的半径分别为$ct'$和$ct$，因此可以推定的是，在$S'$和$S$系中闪光波前球面的方程分别为
\begin{Equation}[光速不变原理]
    c^2t'^2-x'^2-y'^2-z'^2=0\qquad
    c^2t^2-x^2-y^2-z^2=0
\end{Equation}
这里\xref{eq:光速不变原理}给出的两个球面方程，就可以视为对光速不变原理的一种数学表述。\cite{b5}


